\chapter*{先把问题讲明白，再进入证明}
\addcontentsline{toc}{chapter}{先把问题讲明白，再进入证明}
这份讲义从一个很具体的问题出发：一排数，每次选几个，把它们都换成平均值，能不能经过有限次操作使它们完全相等？你只需要先知道平均数、最大公因数、模运算和二阶矩阵乘法。后面会用到群和环，但我们会先说明这些对象解决什么问题，再给定义。

读数学证明时，最容易迷路的地方不是某个计算太长，而是不知道为什么突然要算它。例如，为什么研究平均数的人开始写 $2\times2$ 矩阵？为什么需要一个局部化环？为什么把两个数变成模 $q$ 相同，竟然就离精确相同不远了？本版把这些转折作为讲解的重点。

\section*{先记住两个动作}
第一个动作是\textbf{下降}。如果每一步都让一个非负整数变小，就不能一直做下去。平方和正是我们的第一个下降量。但为了使平方和一直为整数，选择平均组时要满足整除条件，还要留下某些同余差异。

第二个动作是\textbf{输送}。位置较少时，强求每一步都是整数平均会卡住。我们改为把整排数整理成三个等值块，只用两个参数描述它们。几次平均合在一起，就成为两个参数上的一个矩阵。取得足够的矩阵后，直接把参数送到一个能收尾的地方。

这两种方法都要先回答同一个问题：哪些输入从一开始就绝无可能成功？答案来自一个最大公因数 $G$。前两讲会算清楚它，不要求你先接受任何算术群结论。

\section*{第一次和第二次读，任务可以不同}
第一次沿着每讲开头的例子和问题走，目标是能用自己的话解释每个工具为什么出现。第5讲的五平均十五元计数表、第8讲的具体输送矩阵、第10讲的五平均十一元分组，尤其值得亲手算一遍。

第二次再逐行核对各个正式命题的证明。所有原讲义的正式命题、证明段和有限计数表都保留在相应讲次中，讲解插在它们前后。新讲解会把长公式拆成可跟踪的任务，但不会把有限分类替换成随机实验，也不会省略原来的参数分支。

第10讲以后是最费计数的部分。此时先问三件事：这张表从各块取了几个位置？一共取够了没有、取多了没有？取完以后想得到哪一个矩阵？理解这三件事后，再看配比公式，公式就不再是一串没有来由的字母。

\section*{几个会反复出现的词}
\begin{longtable}{p{.20\textwidth}p{.71\textwidth}}
\toprule 讲义用语 & 先用普通语言理解\\\midrule
本原 & 将整数坐标共同的最大公因数约掉。\\
危险素数 & $n$ 的素因子中不等于平均元数 $p$ 的那些。\\
见证 & 两个余数不同的位置，它们证明“现在还没有全同余”。\\
重值 & 某个数值至少出现 $p$ 次。数值的大小不限。\\
块 & 一组当前数值相等的位置。位置并未被合并成一个。\\
返回、宏 & 一段真实平均操作，做完后恢复同一种块重数。\\
射影 & 暂时只看两个参数的比例；实际配置仍带着共同尺度。\\
单位 & 在所讨论的环里有乘法逆元的数。\\
剪切、根 & 给一个坐标加上另一个坐标的某个倍数。\\
同余层 $q$ & 要求矩阵模 $q$ 为单位矩阵；$q$ 越大通常要求越严。\\
同余纤维 & 所有模 $q$ 给出同一个列向量的本原参数。\\
内容 & 整数列两个坐标的最大公因数。\\\bottomrule
\end{longtable}
“安全”在本讲义中是数学定义：操作后仍保留所要求的同余差异。“深核”在原稿中指较深的主同余子群；正文会尽量直接写“已经能实现的 $\Gamma(q,R)$”，让你看见实际用了哪个模数。

这是一份对既有证明的教学改写。Morris 和 Serre 的结果作为外部定理引用，其作用和前提在第8讲说明。小参数中的精确有限证书完整保留。新增例子和教学叙述的检查，不等于对整项研究作了一次独立同行审稿。
